\documentclass{article}

\usepackage[a4 paper, margin=2cm]{geometry}
\usepackage[utf8]{inputenc}
\usepackage{amssymb}
\usepackage{amsmath}
\usepackage{amsthm}
\usepackage{array}
\usepackage{multirow}
\usepackage{multicol}
\usepackage{graphicx}
\usepackage[spanish]{babel}
\usepackage{amsfonts}
\usepackage{float}
\usepackage{mathtools}
\usepackage{scrextend}
\usepackage{bm}
\usepackage{diagbox}
\usepackage{dsfont}
\usepackage{xspace}
\usepackage{booktabs,caption}
\usepackage{tcolorbox}
\usepackage{etoolbox}
\usepackage{adjustbox}
\usepackage{pifont}
\usepackage{framed}
\usepackage{tikz}
\usepackage{tikz-network}
\usepackage{listings}
\newcommand\pythonstyle{\lstset{
    language=Python,
    otherkeywords={self},
    showstringspaces=false,
    backgroundcolor=\color{gray!20},
}}
\lstnewenvironment{python}[1][]{
    \pythonstyle
    \lstset{#1}
}{}
\usetikzlibrary{positioning}
\usepackage{hyperref}
\hypersetup{
    colorlinks=true,
    linkcolor=blue,
    filecolor=magenta,
    urlcolor=blue}

\definecolor{shadecolor}{gray}{0.8}

\setlength{\parskip}{3pt}
\setlength{\parindent}{0pt}
\newcommand{\cpiv}[1]{\colorbox{red!50!white}{#1}}
\newcommand{\rpiv}[1]{\colorbox{blue!50!white}{#1}}
\newcommand{\npiv}[1]{\colorbox{purple!65!white}{#1}}

\title{Trabajo Práctico de Optimización no Lineal \\
\large Introducción a Investigación Operativa y Optimización - Segundo Cuatrimestre 2025}
\author{}
\date{}

\begin{document}

    \maketitle

    La fecha límite de entrega es el 05/12 a las 23:59. En el caso de haber errores graves,
        podrán contar con una fecha de reentrega. Se recomienda utilizar como base el documento de Colab que se
        encuentra en este
        \href{https://colab.research.google.com/drive/1c-Ume72H4fAdg0eKE0YiovByL9ZI_Dxh?usp=sharing}{\textbf{link}}.

    \vspace*{\baselineskip}

    \subsection*{Regresión Lasso}

    En esta sección utilizaremos el dataset \texttt{transfermarkt\_fbref\_201920.csv} con datos sobre gran variedad
    de atributos de jugadores de fútbol, para la temporada 2019-2020. El objetivo será implementar un modelo de
    regresión lineal para predecir el valor de un jugador a partir de sus atributos.

    \textbf{Ejercicio 1.}
    Regresión Lineal Ordinaria.
    \begin{enumerate}
        \item Elegir dos métodos entre: método de gradiente con mitad de paso óptimo, gradiente conjugado o ADAM.
        Utilizarlos para hallar $\beta$ y comparar su desempeño en relación a cantidad de iteraciones y tiempo de
        ejecución. Utilizar \texttt{k\_max=5000} o inferior.
        \item Para $\beta$ hallado en el punto anterior, calcular el cociente de determinación ($R^2$) para los datos
        de entrenamiento y de testeo. ¿Qué se observa? ¿Cómo lo interpretarían?
        \item Interpretar $\beta$: ¿qué atributos benefician el valor del jugador? ¿Qué atributos lo perjudican?
        ¿Cuánta variabilidad hay entre las coordenadas de $\beta$?. Se pueden utilizar las
        funciones \texttt{etiquetas\_beta} y \texttt{graficar\_betas} de la plantilla del Colab.
    \end{enumerate}

    \textbf{Ejercicio 2.} Regresión Lasso.

    \begin{enumerate}
        \item Primera adaptación de ADAM a Lasso: implementar una función \texttt{adam\_penalizado} que aplique el
        método ADAM general, siguiendo el pseudocódigo visto en clase,
        añadiendo un argumento obligatorio \texttt{lamda} para el valor de regularización $\lambda$ y que, en cada
        iteración, sume a \texttt{d} el subgradiente de $\lambda\lVert\beta\rVert_1$.
        \item Aplicar \texttt{adam\_penalizado} a la función:
        \[f(\beta) = \sum_{i=1}^N \left(y^{(i)} - \sum_{j=1}^M\beta_k x_k^{(i)}\right)^2\]
        con $\lambda=10$, $\alpha=0.01$ y $2000$ iteraciones. ¿Es satisfactorio el tiempo de cómputo?
        \item Segunda adaptación de ADAM a Lasso: identificar un cuello de botella en el desempeño de
        \texttt{adam\_penalizado}. Implementar \texttt{lasso\_adam}, que mejore el desempeño. Probar
        \texttt{lasso\_adam} con $\lambda=10$, $\alpha=0.01$ y $2000$ iteraciones. Comparar el tiempo de cómputo con
        la versión del item anterior. \\
        \textbf{Pista:} ¿podríamos aprovechar que nos interesa minimizar una función en particular?
        \item Para el $\beta$ obtenido usando \texttt{lasso\_adam} con $\lambda=10$, calcular el cociente de determinación
        ($R^2$) para los datos de entrenamiento y de testeo. ¿Qué diferencia se observa respecto a Regresión Lineal?
        \item Interpretar el $\beta$ obtenido usando \texttt{lasso\_adam} con $\lambda=10$. ¿Qué atributos benefician y
        perjudican el valor de un defensor? ¿Cuánta variabilidad hay entre las coordenadas de $\beta$? Comparar con la
        interpretación del $\beta$ obtenido con Regresión Lineal. En cada método, ¿cuántos atributos tienen un peso
        con valor absoluto mayor que $0.05$? (para esto se puede usar la función \texttt{etiquetas\_beta\_umbral}).
        \item Implementar la función \texttt{lasso\_coordenadas} que aplique minimización por coordenadas para hallar
        $\beta$ de la Regresión Lasso, con $\lambda=10$. Comparar la interpretación de $\beta$ con la obtenida en el
        ejercicio anterior y el coeficiente de determinación para el conjunto de testeo.
    \end{enumerate}

    \textbf{Ejercicio 3.} Realizar una Regresión Lasso para predecir el valor de los jugadores de otra(s) posición(es).
    Para esto, utilizar \texttt{lasso\_adam} o \texttt{lasso\_coordenadas}.
    Observar qué sucede para distintos valores de $\lambda$ y, a partir de los $\beta$ hallados,
    interpretar cuáles son las características que más influyen en el valor de los jugadores de esa posición.

    \subsection*{Máquina de soporte vectorial (\textit{Support Vector Machine})}

    En esta sección, el objectivo es implementar SVM  para clasificación, adaptando algunos de los métodos de
    optimización que hemos visto en la cursada.
    Utilizaremos el dataset \texttt{fifa\_players.csv} que contiene los atributos de jugadores de
    fútbol, extraídos del FIFA 18.

    \textbf{Ejercicio 4.} Clasificando la posición de los jugadores.

    \begin{enumerate}
        \item Implementar \texttt{subgradiente\_svm}, que calcule el subgradiente de:
        \[f(w,b) = \lambda\frac{1}{2}\lVert w\rVert^2 + \frac{1}{n}\sum_{i=1}^n \max(0, 1-y_i(wx_i+b))\]
        tomando como argumentos la matriz de atributos $X$, el vector de target $y$, $w$, $b$ y el parámetro de
        redularización $\lambda$.
        \item Implementar \texttt{metodo\_gradiente\_svm} y \texttt{adam\_svm} que apliquen el Método del Gradiente
        con paso fijo y ADAM, respectivamente, para hallar los valores de $w$ y $b$. El Método del Gradiente con
        paso fijo debe tomar como argumento la longitud del paso $\alpha$, con valor por defecto $0.01$.
        \item Aplicar ambos métodos para clasificar arqueros y delanteros en base de los atributos \textit{dribbling} y
        \textit{penalties}, utilizando \texttt{k\_max} por lo menos $500$. Para cada uno, visualizar el resultado del
        método y calcular la precisión de la clasificación para los conjuntos de entrenamiento y testeo. Probar con
        distintos valores para los parámetros. ¿ADAM converge siempre a la solución esperada? ¿Por qué creen que sucede esto?
        \item Elegir otras dos posiciones\footnote{En este \href{https://www.fifplay.com/encyclopedia/position/}{link}
        pueden encontrar las definiciones de las posiciones} y otro par de atributos y repetir el ítem anterior.
        Para filtrar las posiciones y los atributos elegidos, pueden utilizar la función
        \texttt{filtrar\_posicion\_atributos}. Para visualizar los atributos
        elegidos, pueden usar \texttt{graficar}. No es necesario que los datos sean estrictamente linealmente
        separables. Pueden probar con distintas elecciones de posiciones y atributos.
    \end{enumerate}

\end{document}

